\chapter{Технологическая часть}

\section{Выбор языка и среды программирования}

Для разработке ПО использовался язык программирования C.~\cite{c} Выбор языка С обусловлен тем, что в настоящий момент большая часть исходного кода ядра Linux, его модулей и драйверов написана на данном языке.~\cite{rust_in_linux}

В качестве компилятора был выбран gcc.~\cite{gcc} Выбор обоснован тем, что данный компилятор является предпочтительным для сборки ядра Linux.~\cite{build_linux}

\section{Реализация алгоритмов определения и сохранения приоритетов, времени непрерывной работы, простоя и блокировки процессов}

В листинге \ref{lst:kernel} представлена реализация алгоритмов определения и сохранения приоритетов, времени непрерывной работы, простоя и блокировки процессов.

\begin{lstlisting}[
	caption={Реализация алгоритмов определения и сохранения приоритетов, времени непрерывной работы, простоя и блокировки процессов},
	label=lst:kernel,
	]
static int printTasks(void *arg)
{
	struct task_struct *task;
	size_t currentPrint = 1;
	char currentString[TMP_STRLEN];
	
	while (currentPrint <= LOGS_COUNT)
	{
		task = &init_task;
		memset(currentString, 0, TMP_STRLEN);
		snprintf(currentString, TMP_STRLEN, "~~~~~~~~~~~~~~~~~~~~~~~~~~%lu "
		"TIME~~~~~~~~~~~~~~~~~~~~~~~~~~\n", currentPrint);
		checkOverflow(currentString, log, LOG_SIZE);
		strcat(log, currentString);
		
		for_each_process(task)
		{
			if (task_is_realtime(task))
			{
				memset(currentString, 0, TMP_STRLEN);
				snprintf(currentString, TMP_STRLEN,
					"pid: %-5d, comm: %15s\nprio: %3d, static_prio: %3d, 
					normal_prio (with "
					"scheduler policy): %3d, rt_priority: %3d\n"
					"run_delay: %10lld\nutime: %10lld (ticks), stime: 
					%15lld (ticks)\n"
					"sched_statistics:\nwait_max: %10llu, wait_count: 
					%10llu, wait_sum: %10llu\n"
					"iowait_count: %10llu, iowait_sum: %10llu\n"
					"sleep_max: %10llu, sum_sleep_runtime: %10llu\n"
					"block_max: %10llu, sum_block_runtime: %10llu, exec_max: 
					%10llu\n\n",
					task->pid, task->comm, task->prio, task->static_prio, 
					task->normal_prio, task->rt_priority,
					task->sched_info.run_delay, task->utime, task->stime,
					task->stats.wait_max, task->stats.wait_count, 
					task->stats.wait_sum,
					task->stats.iowait_count, task->stats.iowait_sum,
					task->stats.sleep_max, task->stats.sum_sleep_runtime,
					task->stats.block_max, task->stats.sum_block_runtime, 
					task->stats.exec_max
				);
				checkOverflow(currentString, log, LOG_SIZE);
				strcat(log, currentString);
			}
		}
		
		currentPrint++;
		mdelay(DELAY_MS);
	}
	
	return 0;
}

static int analyzerOpen(struct inode *spInode, struct file *spFile)
{
	printk(KERN_INFO "%s open called\n", LOG_PREFIX);
	try_module_get(THIS_MODULE);
	return 0;
}

static ssize_t analyzerRead(struct file *filep, char __user *buf, 
	size_t count, loff_t *offp)
{
	ssize_t logLen = strlen(log);
	printk(KERN_INFO "%s read called\n", LOG_PREFIX);
	
	if (copy_to_user(buf, log, logLen))
	{
		printk(KERN_ERR "%s copy_to_user error\n", LOG_PREFIX);
		
		return -EFAULT;
	}
	
	memset(log, 0, LOG_SIZE);
	
	return logLen;
}

static ssize_t analyzerWrite(struct file *file, const char __user *buf, 
	size_t len, loff_t *offp)
{
	printk(KERN_INFO "%s write called\n", LOG_PREFIX);
	return 0;
}

static int analyzerRelease(struct inode *spInode, struct file *spFile)
{
	printk(KERN_INFO "%s release called\n", LOG_PREFIX);
	module_put(THIS_MODULE);
	return 0;
}

static struct proc_ops fops = {
	proc_read:    analyzerRead,
	proc_write:   analyzerWrite,
	proc_open:    analyzerOpen,
	proc_release: analyzerRelease
};
\end{lstlisting}

\section{Реализация алгоритма предоставления информации о процессах пользователю}

В листинге \ref{lst:user} представлена реализация алгоритма предоставления информации об изменении приоритетов, времени непрерывной работы, простоя и блокировки процессов пользователю.

\begin{lstlisting}[
	caption={Реализация алгоритма предоставления информации об изменении приоритетов, времени непрерывной работы, простоя и блокировки процессов пользователю},
	label=lst:user,
	]
int main(int argc, char *argv[])
{
	FILE *filePtr = NULL;
	char log[LOG_SIZE] = {'\0'};
	
	system("sudo insmod taskInfoGetter.ko");
	
	for (int i = 0; i < LOGS_COUNT; i++)
	{
		filePtr = popen("cat /proc/processAnalyzer", "r");
		
		if (filePtr == NULL)
		{
			printf("Error: can't execute cat for process analyzer");
			return 1;
		}
		
		while (fgets(log, sizeof(log), filePtr) != NULL)
		printf("%s", log);
		
		pclose(filePtr);
		sleep(DELAY_S);
	}
	
	system("sudo rmmod taskInfoGetter");
	
	return EXIT_SUCCESS;
}
\end{lstlisting}



\section{Makefile}

В листинге \ref{lst:makefile} представлен Makefile для сборки компонентов программы.

\begin{lstlisting}[
	caption={Содержание Makefile},
	label=lst:makefile,
	]
ifneq ($(KERNELRELEASE),)
	obj-m := taskInfoGetter.o
	taskInfoGetter-objs := md.o
else
	CURRENT = $(shell uname -r)
	KDIR = /lib/modules/$(CURRENT)/build
	PWD = $(shell pwd)

default:
	$(MAKE) -C $(KDIR) M=$(PWD) modules
	make clean

logger:
	gcc starterLogger.c -Wall -Werror -o sl.exe

clean:
	@rm -f *.o .*.cmd .*.flags *.mod.c *.order *.mod *.symvers
	@rm -f .*.*.cmd *~ *.*~ TODO.*
	@rm -fR .tmp*
	@rm -rf .tmp_versions

disclean: clean
	@rm *.ko *.symvers *.mod
endif 
\end{lstlisting}



\section{Демонстрация работы}

В листинге \ref{lst:res} представлена демонстрация работы программы загрузки модуля и получения информации.

\begin{lstlisting}[
caption={Фрагмент вывода при запуске программы},
label=lst:res,
]
~~~~~~~~~~~~~~~~~~~~~~~~~~1 TIME~~~~~~~~~~~~~~~~~~~~~~~~~~
pid: 16   , comm:     migration/0
prio:   0, static_prio: 120, normal_prio (with scheduler policy):   0, 
rt_priority:  99
run_delay: 1112861120
utime:          0 (ticks), stime:         8000000 (ticks)
sched_statistics:
wait_max:      8400, wait_count:        410, wait_sum:         3657080
iowait_count:          0, iowait_sum:          0
sleep_max:       3000, sum_sleep_runtime: 1101204040
block_max:       8400, sum_block_runtime: 1112861120, exec_max:      74000
\end{lstlisting}

В листинге \ref{lst:syslog} представлено содержание журнала ядра после загрузки модуля в ядро.

\begin{lstlisting}[
caption={Содержание журнала ядра},
label=lst:syslog,
]
[  465.991558] [PROCESS ANALYZER] module loaded
[  465.996553] [PROCESS ANALYZER] open called
[  465.996591] [PROCESS ANALYZER] read called
[  465.996660] [PROCESS ANALYZER] read called
[  465.996807] [PROCESS ANALYZER] release called
[  476.001023] [PROCESS ANALYZER] open called
[  476.001075] [PROCESS ANALYZER] read called
[  476.001194] [PROCESS ANALYZER] release called
[  486.005731] [PROCESS ANALYZER] open called
[  486.005777] [PROCESS ANALYZER] read called
[  486.005888] [PROCESS ANALYZER] release called
[  496.010964] [PROCESS ANALYZER] open called
[  496.011017] [PROCESS ANALYZER] read called
[  496.011132] [PROCESS ANALYZER] release called
[  506.016397] [PROCESS ANALYZER] open called
[  506.016450] [PROCESS ANALYZER] read called
[  506.016569] [PROCESS ANALYZER] release called
[  516.021766] [PROCESS ANALYZER] open called
[  516.021817] [PROCESS ANALYZER] read called
[  516.021934] [PROCESS ANALYZER] release called
[  526.051060] [PROCESS ANALYZER] module unloaded
\end{lstlisting}


